\documentclass[10pt]{article}
\usepackage[T1]{fontenc}
\usepackage[utf8]{inputenc}
\usepackage[french]{babel}
\usepackage{amsmath,amssymb,mathrsfs,mathpazo}
\usepackage{geometry}
\geometry{a4paper,margin=2cm}
\usepackage{graphicx,array,tabularx,tabularray,colortbl,multirow,multicol,booktabs,enumitem,pifont,ulem,fancybox,diagbox,cancel}
\usepackage{pstricks,pstricks-add,pst-plot,pst-tree,pst-node,pst-text,pst-func,pst-math,pst-bspline,pst-3dplot,pst-eucl,pst-solides3d}
\let\pscalculate\undefined
\usepackage{pst-bezier}
\usepackage{tikz,pgfplots,tkz-tab}
\usetikzlibrary{arrows,arrows.meta,patterns,calc,positioning,trees,decorations.markings,shapes.geometric,shapes.misc,shapes.arrows,intersections,decorations.pathmorphing,decorations.pathreplacing,backgrounds,shadows,babel}
\pgfplotsset{compat=1.18}
\usepgfplotslibrary{fillbetween}
\usepackage[np]{numprint}
\usepackage{listings,makecell,pas-tableur,twoopt}
\let\barre\undefined
\usepackage[pstricks]{tablvar}
\definecolor{darkgreen}{RGB}{0,100,0}
\providecommand{\red}{\color{red}}
\providecommand{\blue}{\color{blue}}
\providecommand{\green}{\color{green}}
\providecommand{\black}{\color{black}}
\providecommand{\white}{\color{white}}
\providecommand{\gray}{\color{gray}}
\providecommand{\cyan}{\color{cyan}}
\providecommand{\magenta}{\color{magenta}}
\providecommand{\yellow}{\color{yellow}}
\providecommand{\violet}{\color{violet}}
\providecommand{\orange}{\color{orange}}
\providecommand{\euro}{\texteuro}
\providecommand{\EUR}{\texteuro}
\providecommand{\ieme}{\textsuperscript{e}}
\providecommand{\degres}{\ensuremath{^\circ}}
\providecommand{\degree}{\ensuremath{^\circ}}
\providecommand{\vv}[1]{\overrightarrow{#1}}
\providecommand{\Coord}[1]{#1}
\providecommand{\decosix}{}
\providecommand{\starredbullet}{\ensuremath{\bullet}}
\pagestyle{empty}
\setlength{\parindent}{0pt}
\usepackage{adjustbox}\DeclareUnicodeCharacter{2236}{:}\DeclareUnicodeCharacter{2190}{\ensuremath{\leftarrow}}
\providecommand{\R}{\mathbb{R}}
\providecommand{\N}{\mathbb{N}}
\providecommand{\Z}{\mathbb{Z}}
\providecommand{\Q}{\mathbb{Q}}
\providecommand{\D}{\mathbb{D}}
\providecommand{\C}{\mathbb{C}}
\providecommand{\e}{{\mathrm{e}}}
\providecommand{\vect}[1]{\overrightarrow{#1}}
\providecommand{\vectt}[1]{\overrightarrow{\mathrm{#1}}}
\providecommand{\vvt}[1]{\overrightarrow{\mathrm{#1}}}
\providecommand{\np}[2][]{{#2~#1}}
\providecommand{\vv}[1]{{\overrightarrow{#1}}}
\renewcommand{\d}{{\mathrm{d}}}
\providecommand{\degres}{{^\circ}}
\providecommand{\degre}{{^\circ}}
\providecommand{\blue}{\color{blue}}
\providecommand{\red}{\color{red}}
\providecommand{\enskip}{\quad}
\providecommand{\llbracket}{\lbrack\!\lbrack}
\providecommand{\rrbracket}{\rbrack\!\rbrack}
\providecommand{\euro}{{\text{€}}}
\providecommand{\no}{{\text{n°}}}
\providecommand{\interval}[2]{{[#1;#2]}}
\providecommand{\intervFO}[2]{{[#1;#2[}}
\providecommand{\dotfill}{\ldots}

\providecommand{\ds}{\displaystyle}
\providecommand{\textding}[1]{#1}
\providecommand{\Oij}{\(\mathrm{O};\vec{\imath},\vec{\jmath}\)}
\providecommand{\Oijk}{\(\mathrm{O};\vec{\imath},\vec{\jmath},\vec{k}\)}
\providecommand{\Ouv}{\(\mathrm{O};\vec{u},\vec{v}\)}
\providecommand{\intervFF}[2]{\ensuremath{[#1;#2]}}
\renewcommand{\euro}{\ifmmode\text{\texteuro}\else\texteuro\fi}
\DeclareUnicodeCharacter{00B0}{\ifmmode^{\circ}\else\textdegree\fi}
\DeclareUnicodeCharacter{20AC}{\ifmmode\text{\texteuro}\else\texteuro\fi}
% Énoncés : Éducation nationale. Transcription et corrigés : APMEP, Denis Vergès et contributeurs. Découpage et classement : CoopMaths, Sébastien Lozano, Éric Elter et contributeurs. Adaptation de la présentation : Maths Explorer.
% CC BY-SA : https://coopmaths.fr/www/about/
% https://coopmaths.fr/alea/static/bac/2026/tex/bac_2026_06_sujet1_etrangers_2_cor.tex
\begin{document}
\noindent\textit{APMEP / CoopMaths. CC BY-SA. Présentation adaptée par Maths Explorer.}\par\medskip
\begin{enumerate}[itemsep=5mm]
	\item On a
	\[
	f(x)=4\e^{-x}+\cos x+\sin x.
	\]
	Donc :
	\[
	f'(x)=-4\e^{-x}-\sin x+\cos x.
	\]
	Ainsi :
	\[
	f'(x)+f(x)=\left(-4\e^{-x}-\sin x+\cos x\right)+\left(4\e^{-x}+\cos x+\sin x\right)=2\cos x.
	\]
	La fonction $f$ est bien solution de l'équation différentielle $(E)$.

	\medskip
	\textbf{Affirmation 1 : vraie.}

	\item Les points d'intersection des courbes vérifient :
	\[
	2x=\sin x.
	\]
	Posons $h(x)=2x-\sin x$.
	Alors :
	\[
	h'(x)=2-\cos x.
	\]
	Or, pour tout réel $x$, $\cos x\leqslant 1$, donc :
	\[
	h'(x)=2-\cos x\geqslant 1>0.
	\]
	La fonction $h$ est donc strictement croissante sur $\R$.
	De plus :
	\[
	h(0)=0.
	\]
	L'équation $h(x)=0$ admet donc une unique solution, à savoir $x=0$.
	Les deux courbes ont donc un seul point d'intersection, le point de coordonnées $(0~;~0)$.

	\medskip
	\textbf{Affirmation 2 : vraie.}

	\item Pour tout entier naturel $n$ :
	\[
	v_n=\dfrac{2n+\sin n}{n+1}.
	\]
	On écrit :
	\[
	v_n=\dfrac{n\left(2+\frac{\sin n}{n}\right)}{n\left(1+\frac{1}{n}\right)}
	=\dfrac{2+\frac{\sin n}{n}}{1+\frac{1}{n}}.
	\]
	Comme $-1\leqslant \sin n\leqslant 1$, on a
	\[
	\lim\limits_{n\to +\infty}\dfrac{\sin n}{n}=0.
	\]
	Donc :
	\[
	\lim\limits_{n\to +\infty} v_n=\dfrac{2+0}{1+0}=2.
	\]
	La suite $(v_n)$ converge vers 2.

	\medskip
	\textbf{Affirmation 3 : fausse.}

	\item Montrons par récurrence que, pour tout entier naturel $n\geqslant 1$, $u_n=n^2$.

	\emph{Initialisation :} pour $n=1$, on a $u_1=1$ et $1^2=1$. La propriété est vraie au rang 1.

	\emph{Hérédité :} supposons que, pour un entier $n\geqslant 1$, on ait $u_n=n^2$.
	Alors :
	\[
	u_{n+1}=u_n+2n+1=n^2+2n+1=(n+1)^2.
	\]
	La propriété est donc héréditaire.

	Par récurrence, pour tout entier naturel $n\geqslant 1$, on a $u_n=n^2$.

	\medskip
	\textbf{Affirmation 4 : vraie.}

	\item La suite $(u_n)$ est géométrique de premier terme $u_0=1$ et de raison $\e^{-1}$.
	Comme $0<\e^{-1}<1$, on a :
	\[
	\lim\limits_{n\to +\infty} S_n=\dfrac{u_0}{1-\e^{-1}}
	=\dfrac{1}{1-\frac{1}{\e}}
	=\dfrac{\e}{\e-1}.
	\]

	\medskip
	\textbf{Affirmation 5 : vraie.}
\end{enumerate}

\end{document}
